\documentclass[10pt,a4paper]{amsart}
\usepackage[margin=19mm]{geometry}
\usepackage{amsmath,amssymb,mathtools}
\usepackage[scr=rsfs,cal=euler]{mathalfa}
\usepackage{xfrac}
\usepackage[unicode,hidelinks,bookmarks=false]{hyperref}
\newcommand{\ie}{i.e.\ }
\newcommand{\cf}{cf.\ }
\newcommand{\resp}{respectively\ }
\newcommand{\Subsection}{\subsection}
\providecommand{\nobreakdash}{\nobreak\hbox{-}}
\setlength{\emergencystretch}{2em}
\multlinegap0pt
\usepackage{amssymb}
\usepackage{xcolor}
\usepackage{bm}
\usepackage{tikz}
\newtheorem{definition}{Definition}
\newtheorem{lemma}{Lemma}
\title{Combinatorial Capacity Bounds for the $q$-ary Deletion Channel $^{\star}$}
\author{Hassan Tavakoli, Thinh Nguyen, Bella Bose}
\date{Source: arXiv:2607.19559v2 (CC BY 4.0)}
\begin{document}
\maketitle

\noindent\emph{Source and licence.} Combinatorial Capacity Bounds for the $q$-ary Deletion Channel $^{\star}$. Version arXiv:2607.19559v2 (CC BY 4.0), distributed by arXiv under the Creative Commons Attribution 4.0 International licence, http://creativecommons.org/licenses/by/4.0/, as recorded in the arXiv licence metadata for this exact version and shown on the abstract page https://arxiv.org/abs/2607.19559v2. Full licence notice: \texttt{licenses/arxiv-2607.19559-CC-BY-4.0.txt}.\par\smallskip

\noindent\textit{Editorial scope.} Counted unit: the article's lemma lem:Phi\_rec (source0.tex lines 608--615). The article's complete original proof of it is delivered with it (source0.tex lines 617--685, one \texttt{proof} environment of 69 lines). No mathematics is written, repaired or re-derived: the statement and the proof are the article's own lines, in the article's own notation, and no retained line is substituted or re-flowed.  Retained uncounted as the source-local context the proof needs: lines 210--213; lines 427--444.  Nothing was omitted from any retained span, so the package carries no omission record.  No retained line contains a citation, so the wrapper carries no bibliography and no bibliography entry.  No retained line contains a figure environment, an image-inclusion command or drawing code, so no image is delivered and the package declares no asset dependency.  Editorial formatting only, in addition: an AMS \texttt{amsart} preamble that restates the article's own declarations and macros used by the retained lines, namely the environments \texttt{definition}, \texttt{lemma} on separate counters, together with the standard packages those lines need.  The retained environments print in this document as Definition 1, Lemma 1 in order of appearance, whereas the article prints the same items with the numbering of its own full document.  The complete native source of the article as distributed in the arXiv e-print for this exact version is bundled unchanged as \texttt{sources/205747/source0.tex}.
\begin{equation}
  N_n(sx',ty')=N_{n-1}(x',ty')+\mathbf{1}_{\{s=t\}}N_{n-1}(x',y').
  \label{eq:rec}
\end{equation}

\begin{definition}[Entropy correction]
\label{def:Phi}
For $0\le k\le n$ and $x\in\Sigma_q^n$, define the
\emph{per-input correction}
\begin{equation}
  \phi_{k,n}(x)\triangleq \frac{1}{\binom{n}{k}}
  \sum_{y\in\Sigma_q^k}N_n(x,y)\log_2 N_n(x,y)\ge 0,
  \label{eq:phi}
\end{equation}
and the \emph{average correction}
\begin{equation}
  \Phi_{k,n}\triangleq \frac{1}{q^n\binom{n}{k}}
  \sum_{x\in\Sigma_q^n}\sum_{y\in\Sigma_q^k}
  N_n(x,y)\log_2 N_n(x,y)\ge 0.
  \label{eq:Phi}
\end{equation}
Set $\Delta_n(d)\triangleq \sum_{k=1}^{n-1}w_k\Phi_{k,n}\ge 0$.
\end{definition}

\begin{lemma}[Recursive lower bound on $\Phi_{k,n}$]
\label{lem:Phi_rec}
For $1\le k\le n-1$:
\begin{equation}
  \Phi_{k,n}\ge\frac{n-k}{n}\Phi_{k,n-1}+\frac{k}{n}\Phi_{k-1,n-1}.
  \label{eq:Phi_rec}
\end{equation}
\end{lemma}

\begin{proof}
Define the auxiliary quantity
\(
  S_{k,n}\triangleq\sum_{x\in\Sigma_q^n}\sum_{y\in\Sigma_q^k}
  N_n(x,y)\log_2 N_n(x,y),
\)
so that $\Phi_{k,n}=S_{k,n}/(q^n\binom{n}{k})$ by
Definition~\ref{def:Phi}.  It suffices to show
\begin{equation}
  S_{k,n}\ge q\,S_{k,n-1}+q\,S_{k-1,n-1}.
  \label{eq:Srec}
\end{equation}

\paragraph{Super-additivity of $f(u)=u\log_2 u$.}
We claim $f(a+b)\ge f(a)+f(b)$ for all $a,b\ge 0$.  The cases
$a=0$ or $b=0$ are trivial since $f(0)=0$.  For $a,b>0$:
\(
  f(a+b)-f(a)-f(b)
  =(a+b)\log_2(a+b)-a\log_2 a-b\log_2 b
  =a\log_2\!\Bigl(1+\frac{b}{a}\Bigr)
   +b\log_2\!\Bigl(1+\frac{a}{b}\Bigr)\;\ge\;0,
\)
since both terms are non-negative for $a,b>0$.

\paragraph{Applying super-additivity entrywise.}
From the recursion~\eqref{eq:rec},
for $x=sx'\in\Sigma_q^n$ and $y=ty'\in\Sigma_q^k$.
Super-additivity of $f$ gives, for every $s,t,x',y'$:
\begin{align}
  &f\bigl(N_n(sx',ty')\bigr) \nonumber\\
  &\ge f\bigl(N_{n-1}(x',ty')\bigr)
   +\mathbf{1}_{\{s=t\}}f\bigl(N_{n-1}(x',y')\bigr).
  \label{eq:fineq}
\end{align}

\paragraph{Summing to obtain~\eqref{eq:Srec}.}
Sum~\eqref{eq:fineq} over all $s,t\in\Sigma_q$,
$x'\in\Sigma_q^{n-1}$, $y'\in\Sigma_q^{k-1}$.  The left-hand side
gives $S_{k,n}$.  On the right-hand side:
\begin{itemize}
  \item \emph{First term.} Summing $f(N_{n-1}(x',ty'))$ over
        $s\in\Sigma_q$ contributes a factor of $q$ (since the summand
        is independent of $s$), and then summing over $t\in\Sigma_q$,
        $x'\in\Sigma_q^{n-1}$, $y'\in\Sigma_q^{k-1}$ gives
        $q\,S_{k,n-1}$.
  \item \emph{Second term.} The indicator $\mathbf{1}_{\{s=t\}}$
        contributes exactly one value of $s$ per $t$.  Summing over
        $t\in\Sigma_q$, $x'\in\Sigma_q^{n-1}$, $y'\in\Sigma_q^{k-1}$
        then gives $q\,S_{k-1,n-1}$.
\end{itemize}
Hence $S_{k,n}\ge q\,S_{k,n-1}+q\,S_{k-1,n-1}$.

\paragraph{Converting back to $\Phi$.}
Dividing both sides of~\eqref{eq:Srec} by $q^n\binom{n}{k}$:
\(
  \Phi_{k,n}
  \ge \frac{q\,S_{k,n-1}}{q^n\binom{n}{k}}
      +\frac{q\,S_{k-1,n-1}}{q^n\binom{n}{k}}
  = \frac{q^{n-1}\binom{n-1}{k}\,\Phi_{k,n-1}}{q^{n-1}\binom{n}{k}}
    +\frac{q^{n-1}\binom{n-1}{k-1}\,\Phi_{k-1,n-1}}{q^{n-1}\binom{n}{k}}
  = \frac{\binom{n-1}{k}}{\binom{n}{k}}\,\Phi_{k,n-1}
    +\frac{\binom{n-1}{k-1}}{\binom{n}{k}}\,\Phi_{k-1,n-1}
  = \frac{n-k}{n}\,\Phi_{k,n-1}+\frac{k}{n}\,\Phi_{k-1,n-1},
\)
where we used $S_{k,n-1}=q^{n-1}\binom{n-1}{k}\,\Phi_{k,n-1}$,
$S_{k-1,n-1}=q^{n-1}\binom{n-1}{k-1}\,\Phi_{k-1,n-1}$, and the
binomial ratios $\binom{n-1}{k}/\binom{n}{k}=(n-k)/n$,
$\binom{n-1}{k-1}/\binom{n}{k}=k/n$.
\end{proof}

\end{document}
